% ═══════════════════════════════════════════════════════════════════
% Group Chat Application Report — LaTeX Source
% Tutorial: Group Chat Application using WebSockets
% ═══════════════════════════════════════════════════════════════════

\documentclass[12pt, a4paper]{article}

% ── Packages ──────────────────────────────────────────────────────
\usepackage[utf8]{inputenc}
\usepackage[T1]{fontenc}
\usepackage{lmodern}
\usepackage[margin=1in]{geometry}
\usepackage{graphicx}
\usepackage{hyperref}
\usepackage{xcolor}
\usepackage{listings}
\usepackage{booktabs}
\usepackage{enumitem}
\usepackage{float}
\usepackage{caption}
\usepackage{fancyhdr}
\usepackage{titlesec}
\usepackage{tocloft}
\usepackage{amsmath}
\usepackage{parskip}
\usepackage{tabularx}
\usepackage{longtable}

% ── Colours ───────────────────────────────────────────────────────
\definecolor{codebg}{HTML}{F5F5F5}
\definecolor{codeframe}{HTML}{CCCCCC}
\definecolor{keyword}{HTML}{0077AA}
\definecolor{comment}{HTML}{999999}
\definecolor{string}{HTML}{669900}
\definecolor{linkblue}{HTML}{0066CC}

% ── Hyperref Setup ────────────────────────────────────────────────
\hypersetup{
    colorlinks=true,
    linkcolor=linkblue,
    urlcolor=linkblue,
    citecolor=linkblue,
    pdfauthor={Team Members},
    pdftitle={Group Chat Application using WebSockets — Report},
    pdfsubject={Computer System design Tutorial},
}

% ── Code Listing Style ───────────────────────────────────────────
\lstdefinestyle{codestyle}{
    backgroundcolor=\color{codebg},
    frame=single,
    rulecolor=\color{codeframe},
    basicstyle=\ttfamily\footnotesize,
    keywordstyle=\color{keyword}\bfseries,
    commentstyle=\color{comment}\itshape,
    stringstyle=\color{string},
    breaklines=true,
    breakatwhitespace=false,
    numbers=left,
    numberstyle=\tiny\color{comment},
    numbersep=8pt,
    showstringspaces=false,
    tabsize=4,
    captionpos=b,
}
\lstset{style=codestyle}

% ── Header / Footer ──────────────────────────────────────────────
\pagestyle{fancy}
\fancyhf{}
\fancyhead[L]{\small Group Chat Application — WebSockets}
\fancyhead[R]{\small \thepage}
\fancyfoot[C]{\small Computer System design Tutorial Report}
\renewcommand{\headrulewidth}{0.4pt}
\renewcommand{\footrulewidth}{0.4pt}

% ── Section formatting ───────────────────────────────────────────
\titleformat{\section}{\Large\bfseries}{\thesection.}{0.5em}{}
\titleformat{\subsection}{\large\bfseries}{\thesubsection.}{0.5em}{}
\titleformat{\subsubsection}{\normalsize\bfseries}{\thesubsubsection.}{0.5em}{}

% ═══════════════════════════════════════════════════════════════════
%  DOCUMENT
% ═══════════════════════════════════════════════════════════════════

\begin{document}

% ── Title Page ────────────────────────────────────────────────────
\begin{titlepage}
    \centering
    \vspace*{1cm}

    {\LARGE\bfseries Tutorial Report}\\[0.5cm]
    {\Huge\bfseries Group Chat Application\\using WebSockets}\\[1.5cm]

    \rule{\textwidth}{0.5pt}\\[0.5cm]

    {\large Subject: Computer System design}\\[0.3cm]
    % ── UPDATE the course code and semester below ──
    {\large Course Code: \texttt{XXXX}}\\[0.3cm]
    {\large Semester: \texttt{XXXX}}\\[0.5cm]

    \rule{\textwidth}{0.5pt}\\[1.5cm]

    {\large\bfseries Team Members}\\[0.5cm]

    \begin{tabular}{lll}
        \toprule
        \textbf{Name} & \textbf{Roll No.} & \textbf{Role} \\
        \midrule
        % ── UPDATE names and roll numbers below ──
        Member 1 & XXXXXXXXXX & Group Head / Backend \\
        Member 2 & XXXXXXXXXX & Frontend \\
        Member 3 & XXXXXXXXXX & Testing \\
        Member 4 & XXXXXXXXXX & Documentation \\
        \bottomrule
    \end{tabular}

    \vfill

    % ── UPDATE the institution name below ──
    {\large Institution Name}\\[0.3cm]
    {\large Department of Computer Science \& Engineering}\\[0.3cm]
    {\large August 2026}

\end{titlepage}

% ── Table of Contents ─────────────────────────────────────────────
\newpage
\tableofcontents
\newpage


% ══════════════════════════════════════════════════════════════════
% 1. INTRODUCTION
% ══════════════════════════════════════════════════════════════════
\section{Introduction}

This report presents the design, implementation, and working of a \textbf{real-time Group Chat Application} built using \textbf{WebSockets}. The application enables multiple users to communicate simultaneously through a common chat room with instant message delivery, fulfilling the requirements of the Computer System design tutorial assignment.

The project demonstrates core networking concepts including persistent bidirectional communication (WebSockets), client-server architecture, message broadcasting, and graceful connection management. The application is named \textbf{``PixelChat''} and features a retro 8-bit pixel-art themed user interface, complete with gamification elements such as XP tracking, ranks, and streaks to enhance user engagement.

\subsection{Objectives}

\begin{itemize}[leftmargin=*]
    \item Design and implement a real-time group chat application using WebSockets.
    \item Support \textbf{multiple users communicating simultaneously} through a common chat room.
    \item Set up a single backend server to handle all incoming connections and messages.
    \item Set up front-end clients for 4 users on allotted laboratory machines.
    \item Implement user join/leave notifications, real-time message broadcasting, usernames for identifying participants, and graceful handling of client disconnections.
\end{itemize}

\subsection{Scope}

Beyond the minimum required features, the application also includes:
\begin{itemize}[leftmargin=*]
    \item Avatar selection with 12 predefined character avatars.
    \item Emoji picker with 30 quick-react emojis.
    \item File attachment support (images, videos, audio, documents).
    \item Message delivery receipts (sent, partial, delivered to all).
    \item Typing indicators showing when another user is composing a message.
    \item Message history for users joining mid-conversation.
    \item Auto-reconnect with exponential backoff on connection loss.
    \item Gamification: XP bar, rank progression, message streaks, and session timer.
    \item 8-bit retro sound effects (Mario-themed) for join, leave, and message events.
    \item Connection status indicator (Online / Offline / Reconnecting).
\end{itemize}


% ══════════════════════════════════════════════════════════════════
% 2. TECHNOLOGY STACK
% ══════════════════════════════════════════════════════════════════
\section{Technology Stack}

The application uses a lightweight, modern technology stack requiring only Python and a web browser — no Node.js or heavyweight frameworks.

\begin{table}[H]
    \centering
    \caption{Technology Stack}
    \begin{tabularx}{\textwidth}{l l X}
        \toprule
        \textbf{Layer} & \textbf{Technology} & \textbf{Purpose} \\
        \midrule
        Backend Server  & Python 3 + FastAPI + Uvicorn & Asynchronous WebSocket server with REST endpoint for file uploads \\
        Frontend Server & Python 3 + FastAPI + Uvicorn & Lightweight static file server for HTML/CSS/JS \\
        Frontend Client & HTML5 + CSS3 + Vanilla JavaScript & Responsive retro-themed chat UI with WebSocket client \\
        Protocol        & WebSocket (RFC 6455) & Full-duplex, persistent TCP connection for real-time messaging \\
        Data Format     & JSON & All client-server messages serialized as JSON objects \\
        Configuration   & \texttt{python-dotenv} + \texttt{.env} file & Port configuration via environment variables \\
        \bottomrule
    \end{tabularx}
\end{table}

\subsection{Key Dependencies}

The server requires only four Python packages, listed in \texttt{requirements.txt}:

\begin{table}[H]
    \centering
    \caption{Python Dependencies}
    \begin{tabularx}{\textwidth}{l X}
        \toprule
        \textbf{Package} & \textbf{Purpose} \\
        \midrule
        \texttt{fastapi}          & Modern, high-performance web framework with native WebSocket support \\
        \texttt{uvicorn[standard]} & ASGI server to run FastAPI asynchronously \\
        \texttt{python-dotenv}     & Load \texttt{.env} configuration files \\
        \texttt{python-multipart}  & Required for file upload handling (\texttt{UploadFile}) \\
        \bottomrule
    \end{tabularx}
\end{table}


% ══════════════════════════════════════════════════════════════════
% 3. SYSTEM ARCHITECTURE
% ══════════════════════════════════════════════════════════════════
\section{System Architecture}

The application follows a \textbf{client-server architecture} with a clear separation between the backend WebSocket server and the frontend static file server. Both servers run on the same host machine, but on different ports.

\subsection{Architecture Overview}

\begin{figure}[H]
    \centering
    % ── PLACEHOLDER: Replace with an architecture diagram screenshot ──
    \fbox{\parbox{0.85\textwidth}{\centering\vspace{3cm}
        \textbf{[PLACEHOLDER — Architecture Diagram]}\\[0.5cm]
        Insert a screenshot of the architecture diagram here.\\
        You may use the ASCII diagram from the README or create a visual version.
    \vspace{3cm}}}
    \caption{System Architecture — Client-Server Model}
    \label{fig:architecture}
\end{figure}

The architecture consists of:

\begin{enumerate}[leftmargin=*]
    \item \textbf{Backend Server} (\texttt{server/server.py}) — Runs on port \texttt{8000} (configurable). Handles all WebSocket connections via the \texttt{/ws} endpoint and file uploads via the \texttt{/upload} REST endpoint.

    \item \textbf{Frontend Server} (\texttt{client/serve.py}) — Runs on port \texttt{5000} (configurable). Serves the static HTML, CSS, JavaScript, and sound files to client browsers.

    \item \textbf{Browser Clients} — Each user opens \texttt{http://<server-ip>:5000} in their browser. The JavaScript client establishes a WebSocket connection to \texttt{ws://<server-ip>:8000/ws} for real-time communication.
\end{enumerate}

\subsection{Project Directory Structure}

\begin{lstlisting}[language={}, caption={Project Directory Structure}, label={lst:structure}]
group-chat-app/
|-- .env                        # Port configuration (not committed)
|-- .env.example                # Template for .env
|-- .gitignore                  # Git ignore rules
|-- README.md                   # Project documentation
|-- server/
|   |-- requirements.txt        # Python dependencies
|   |-- server.py               # FastAPI WebSocket server
|   +-- uploads/                # Uploaded files stored here
|-- client/
|   |-- index.html              # Chat UI (login + chat screens)
|   |-- style.css               # 8-bit pixel art themed styles
|   |-- app.js                  # WebSocket client logic + gamification
|   |-- serve.py                # Static file server
|   +-- sounds/                 # Mario-themed 8-bit sound effects
|       |-- mushroom.mp3        # Join sound
|       |-- coin.mp3            # Message sound
|       +-- pipe.mp3            # Leave sound
\end{lstlisting}

\subsection{Communication Flow}

\begin{enumerate}[leftmargin=*]
    \item The user opens the frontend URL in a browser. The browser downloads \texttt{index.html}, \texttt{style.css}, and \texttt{app.js} from the frontend server over HTTP.
    \item The user enters a username, selects an avatar, and clicks ``Start Quest''.
    \item The JavaScript client establishes a WebSocket connection to the backend and sends a \texttt{join} message with the chosen username and avatar.
    \item The backend validates the username (non-empty, unique), registers the user, and broadcasts a join notification to all other connected clients.
    \item When a user sends a chat message, the client transmits it over the WebSocket. The backend broadcasts it to \textbf{all} connected clients (including the sender for receipt confirmation).
    \item When a user disconnects (closes tab, navigates away, or clicks ``Leave''), the backend detects the disconnection, removes the user, and broadcasts a leave notification.
\end{enumerate}


% ══════════════════════════════════════════════════════════════════
% 4. WEBSOCKET PROTOCOL & MESSAGE FORMAT
% ══════════════════════════════════════════════════════════════════
\section{WebSocket Protocol \& Message Format}

All communication between the client and server uses \textbf{JSON-encoded messages} transmitted over a single persistent WebSocket connection. Each message contains a \texttt{type} field that determines how it should be processed.

\subsection{Client $\rightarrow$ Server Messages}

\begin{table}[H]
    \centering
    \caption{Client to Server Message Types}
    \begin{tabularx}{\textwidth}{l l X}
        \toprule
        \textbf{Type} & \textbf{Fields} & \textbf{Description} \\
        \midrule
        \texttt{join}    & \texttt{username}, \texttt{avatar} & Register with a username and chosen avatar. Must be the first message sent after connection. \\
        \texttt{message} & \texttt{text}, \texttt{client\_msg\_id}, \texttt{attachment} & Send a chat message with optional file attachment. \texttt{client\_msg\_id} is a locally generated ID used for delivery receipt tracking. \\
        \texttt{typing}  & — & Signal that the user is currently typing. Throttled to once every 2 seconds on the client side. \\
        \bottomrule
    \end{tabularx}
\end{table}

\subsection{Server $\rightarrow$ Client Messages}

\begin{table}[H]
    \centering
    \caption{Server to Client Message Types}
    \begin{tabularx}{\textwidth}{l l X}
        \toprule
        \textbf{Type} & \textbf{Fields} & \textbf{Description} \\
        \midrule
        \texttt{system}   & \texttt{message}, \texttt{timestamp} & System notification (e.g., welcome message). \\
        \texttt{join}     & \texttt{username}, \texttt{avatar}, \texttt{message}, \texttt{timestamp} & Notification that a user has joined. \\
        \texttt{leave}    & \texttt{username}, \texttt{message}, \texttt{timestamp} & Notification that a user has left. \\
        \texttt{message}  & \texttt{username}, \texttt{avatar}, \texttt{text}, \texttt{attachment}, \texttt{timestamp} & Broadcast of a chat message to all clients. \\
        \texttt{userList} & \texttt{users} (array of \texttt{\{username, avatar\}}) & Updated list of all currently online users. \\
        \texttt{history}  & \texttt{messages} (array) & Last 50 messages, sent to newly joining users. \\
        \texttt{receipt}  & \texttt{msg\_id}, \texttt{status} & Delivery receipt for a sent message. \\
        \texttt{error}    & \texttt{message} & Error notification (e.g., username already taken). \\
        \texttt{typing}   & \texttt{username} & Typing indicator relayed from another user. \\
        \bottomrule
    \end{tabularx}
\end{table}

\subsection{Message Lifecycle Example}

\begin{lstlisting}[language={}, caption={Example: User ``Alice'' sending a message}, label={lst:msgflow}]
1. Client -> Server:
   { "type": "message", "text": "Hello everyone!",
     "client_msg_id": "msg-1723645...-ab12c" }

2. Server -> ALL Clients (broadcast):
   { "type": "message", "username": "Alice", "avatar": "wizard",
     "text": "Hello everyone!", "timestamp": "14:30:05" }

3. Server -> Alice only (receipt):
   { "type": "receipt", "msg_id": "msg-1723645...-ab12c",
     "status": "delivered_all" }
\end{lstlisting}


% ══════════════════════════════════════════════════════════════════
% 5. BACKEND IMPLEMENTATION
% ══════════════════════════════════════════════════════════════════
\section{Backend Implementation}

The backend is implemented in a single Python file (\texttt{server/server.py}, 358 lines) using FastAPI's native WebSocket support.

\subsection{ConnectionManager Class}

The core of the backend is the \texttt{ConnectionManager} class, which maintains the state of all active connections and provides methods for message routing.

\begin{table}[H]
    \centering
    \caption{ConnectionManager Methods}
    \begin{tabularx}{\textwidth}{l X}
        \toprule
        \textbf{Method} & \textbf{Description} \\
        \midrule
        \texttt{add(ws, username, avatar)}   & Register a new WebSocket connection with a username and avatar. \\
        \texttt{remove(ws)}                   & Remove a disconnected user and return their username. \\
        \texttt{get\_username(ws)}            & Retrieve the username associated with a WebSocket. \\
        \texttt{get\_all\_users()}            & Return a sorted list of all online users (username + avatar). \\
        \texttt{is\_username\_taken(name)}    & Case-insensitive check for duplicate usernames. \\
        \texttt{broadcast(msg, exclude)}      & Send a JSON message to all clients, optionally excluding one. Returns the count of successfully delivered messages. \\
        \texttt{send\_to\_all(msg)}           & Send to ALL clients without exclusion. \\
        \texttt{add\_to\_history(msg)}        & Store a message in the rolling history buffer (max 50). \\
        \texttt{start\_cleanup\_timer()}      & Start a 60-second timer to clear history when the room empties. \\
        \bottomrule
    \end{tabularx}
\end{table}

\subsection{WebSocket Endpoint Lifecycle}

The \texttt{/ws} endpoint handles each client through a three-phase lifecycle:

\begin{enumerate}[leftmargin=*]
    \item \textbf{Phase 1 — Join:} Accept the WebSocket connection, receive the first message (must be type \texttt{join}), validate the username (non-empty and unique), and register the user.

    \item \textbf{Phase 2 — Welcome \& Sync:} Send a welcome system message to the joiner, broadcast a join notification to all other users, send the updated user list to everyone, and send message history to the new joiner.

    \item \textbf{Phase 3 — Chat Loop:} Enter an infinite loop listening for incoming messages. For \texttt{message} type, broadcast to all and send a delivery receipt back to the sender. For \texttt{typing} type, relay the typing indicator to all other users.

    \item \textbf{Cleanup (finally block):} On disconnection (normal or error), remove the user, broadcast a leave notification, send the updated user list, and start the history cleanup timer if the room is now empty.
\end{enumerate}

\subsection{File Upload Endpoint}

The server exposes a \texttt{POST /upload} REST endpoint for file attachments:

\begin{itemize}[leftmargin=*]
    \item Accepts any file type via multipart form upload.
    \item Saves the file with a UUID-based unique filename to prevent collisions.
    \item Returns the file URL, original filename, MIME type, and file size as JSON.
    \item Uploaded files are served statically from the \texttt{/uploads/} path.
\end{itemize}

\subsection{Delivery Receipt System}

The server implements a message delivery receipt mechanism:

\begin{table}[H]
    \centering
    \caption{Delivery Receipt Statuses}
    \begin{tabularx}{\textwidth}{l l X}
        \toprule
        \textbf{Status} & \textbf{Emoji} & \textbf{Meaning} \\
        \midrule
        \texttt{sent}           & \raisebox{-0.1em}{\Large 😴} & Message reached the server but no other users are online. \\
        \texttt{partial}        & \raisebox{-0.1em}{\Large 😃} & Message was delivered to some but not all online users. \\
        \texttt{delivered\_all} & \raisebox{-0.1em}{\Large 😎} & Message was successfully delivered to all online users. \\
        \bottomrule
    \end{tabularx}
\end{table}

The server calculates receipt status by comparing the number of successful deliveries against the total number of other connected clients.

\subsection{History Management}

\begin{itemize}[leftmargin=*]
    \item The server maintains a rolling buffer of the last 50 chat messages.
    \item When a new user joins, the entire history is sent as a single \texttt{history} message.
    \item When the last user leaves the room, a 60-second cleanup timer starts.
    \item If no one rejoins within 60 seconds, the history is cleared for a fresh session.
    \item If someone rejoins before the timer expires, the timer is cancelled and the history is preserved.
\end{itemize}


% ══════════════════════════════════════════════════════════════════
% 6. FRONTEND IMPLEMENTATION
% ══════════════════════════════════════════════════════════════════
\section{Frontend Implementation}

The frontend consists of three files — \texttt{index.html} (structure), \texttt{style.css} (styling), and \texttt{app.js} (logic) — served by a lightweight FastAPI static file server (\texttt{serve.py}).

\subsection{User Interface Design}

The UI follows an \textbf{8-bit retro pixel-art aesthetic} using:
\begin{itemize}[leftmargin=*]
    \item \textbf{Font:} ``Press Start 2P'' (Google Fonts) — a pixel-art typeface.
    \item \textbf{Color Palette:} Four-color scheme inspired by retro displays:
    \begin{itemize}
        \item \texttt{\#FDF6ED} — Cream (lightest, text)
        \item \texttt{\#DCCFC0} — Warm ivory/tan
        \item \texttt{\#A1BC98} — Sage green (accent)
        \item \texttt{\#778873} — Dark sage/forest (primary)
    \end{itemize}
    \item \textbf{Visual Effects:} Scanline overlay, blinking cursor animations, pixel-border box shadows, and smooth CSS transitions.
\end{itemize}

The application has two screens:

\subsubsection{Login Screen}

\begin{figure}[H]
    \centering
    % ── PLACEHOLDER: Replace with a login screen screenshot ──
    \fbox{\parbox{0.85\textwidth}{\centering\vspace{4cm}
        \textbf{[PLACEHOLDER — Login Screen Screenshot]}\\[0.5cm]
        Insert a screenshot of the login/join screen here,\\
        showing the title, avatar picker, username input, and Start Quest button.
    \vspace{4cm}}}
    \caption{Login Screen — Avatar Selection \& Username Entry}
    \label{fig:login}
\end{figure}

The login screen features:
\begin{itemize}[leftmargin=*]
    \item A large ``PIXEL CHAT'' title with blinking cursor animation.
    \item A stat row showing live player count, server status, and season.
    \item An \textbf{avatar picker grid} with 12 themed avatars (Wizard, Robot, Ninja, Astronaut, Dragon, Hero, Alien, Cyber, Fox, Owl, Bear, Lion).
    \item A username text input field (max 20 characters).
    \item A ``START QUEST'' button that initiates the WebSocket connection.
    \item Error display area for validation messages (e.g., ``Username already taken'').
\end{itemize}

\subsubsection{Chat Screen}

\begin{figure}[H]
    \centering
    % ── PLACEHOLDER: Replace with a chat screen screenshot ──
    \fbox{\parbox{0.85\textwidth}{\centering\vspace{4cm}
        \textbf{[PLACEHOLDER — Chat Screen Screenshot]}\\[0.5cm]
        Insert a screenshot of the main chat interface here,\\
        showing the sidebar, chat messages, and input area.
    \vspace{4cm}}}
    \caption{Chat Screen — Main Chat Interface}
    \label{fig:chat}
\end{figure}

The chat screen is divided into two panels:

\textbf{Left Sidebar Panel:}
\begin{itemize}[leftmargin=*]
    \item PixelChat logo.
    \item XP progress bar with current rank display (Newbie $\rightarrow$ Squire $\rightarrow$ Knight $\rightarrow$ Champion $\rightarrow$ Warlord $\rightarrow$ Legend).
    \item Online players list with avatars and a ``YOU'' tag for the current user.
    \item Session stats: total messages sent, current streak, and session duration.
\end{itemize}

\textbf{Main Chat Panel:}
\begin{itemize}[leftmargin=*]
    \item Header bar with room name (``\# GREAT HALL''), player count, connection status indicator (Online/Offline/Reconnecting), a receipt info button, and a ``Leave'' button.
    \item Scrollable messages area displaying system notifications, join/leave events, and user messages with avatars and timestamps.
    \item Typing indicator bar (``USERNAME TYPING...'') with animated dots.
    \item Input footer with attachment button, emoji picker toggle, message text input, and send button.
\end{itemize}

\subsection{WebSocket Client (\texttt{app.js})}

The JavaScript client (\texttt{app.js}, 878 lines) handles all real-time communication and UI interactions:

\subsubsection{Connection Management}
\begin{itemize}[leftmargin=*]
    \item \textbf{Connect:} Creates a new \texttt{WebSocket} object targeting \texttt{ws://<hostname>:8000/ws}. On open, sends a \texttt{join} message.
    \item \textbf{Auto-Reconnect:} On unexpected disconnection, schedules reconnection with \textbf{exponential backoff} (1s, 2s, 4s, 8s, max 10s).
    \item \textbf{Intentional Disconnect:} Clicking ``Leave'' sets an \texttt{isIntentionalClose} flag to prevent auto-reconnect.
    \item \textbf{Status Display:} The connection indicator updates in real-time: green dot = Online, red dot = Offline, yellow dot = Reconnecting.
\end{itemize}

\subsubsection{Message Handling}
\begin{itemize}[leftmargin=*]
    \item \textbf{Optimistic Rendering:} Own messages are rendered immediately in the UI before server confirmation, providing instant feedback.
    \item \textbf{Echo Suppression:} When the server broadcasts back the sender's own message, it is suppressed to avoid duplicates.
    \item \textbf{Receipt Tracking:} Each outgoing message is assigned a unique \texttt{client\_msg\_id}. The receipt emoji on the message bubble updates when the server responds.
    \item \textbf{HTML Escaping:} All user-generated content is escaped via \texttt{escapeHtml()} to prevent XSS attacks.
    \item \textbf{History Rendering:} On join, historical messages are rendered with visual separators (``--- PREV MESSAGES ---'' / ``--- NEW MESSAGES ---'').
\end{itemize}

\subsubsection{Typing Indicators}
\begin{itemize}[leftmargin=*]
    \item Typing events are \textbf{throttled} to at most once every 2 seconds to reduce WebSocket traffic.
    \item The indicator auto-hides after 3 seconds of inactivity.
    \item The indicator is also hidden when a message from the typing user arrives.
\end{itemize}

\subsubsection{File Attachments}
\begin{itemize}[leftmargin=*]
    \item Users click the attachment button (📎) to open a file picker.
    \item The selected file is uploaded to the backend via a \texttt{POST /upload} request.
    \item A preview bar shows the filename and upload progress.
    \item The attachment metadata (URL, filename, type, size) is sent alongside the chat message over WebSocket.
    \item Recipients render attachments contextually: inline images, video/audio players, or download cards for other file types.
\end{itemize}

\subsection{Sound Effects}

The application uses three Mario-themed 8-bit sound effects:

\begin{table}[H]
    \centering
    \caption{Sound Effects}
    \begin{tabularx}{\textwidth}{l l X}
        \toprule
        \textbf{Event} & \textbf{Sound File} & \textbf{Description} \\
        \midrule
        User Join   & \texttt{mushroom.mp3} & Super Mario mushroom power-up sound \\
        New Message & \texttt{coin.mp3}     & Mario coin collection sound \\
        User Leave  & \texttt{pipe.mp3}     & Mario Bros warp pipe sound \\
        \bottomrule
    \end{tabularx}
\end{table}

Additionally, synthesized 8-bit sounds are generated via the Web Audio API:
\begin{itemize}[leftmargin=*]
    \item \textbf{Notification beep:} A 440Hz square-wave beep when a message arrives while the tab is not focused.
    \item \textbf{Level-up fanfare:} A four-note ascending arpeggio (C4-E4-G4-C5) played on rank promotion.
\end{itemize}


% ══════════════════════════════════════════════════════════════════
% 7. GAMIFICATION SYSTEM
% ══════════════════════════════════════════════════════════════════
\section{Gamification System}

The application includes a lightweight gamification layer to enhance user engagement, implemented entirely on the client side.

\subsection{XP \& Rank Progression}

Users earn Experience Points (XP) through the following actions:

\begin{table}[H]
    \centering
    \caption{XP Earning Actions}
    \begin{tabularx}{\textwidth}{l c X}
        \toprule
        \textbf{Action} & \textbf{XP Gained} & \textbf{Details} \\
        \midrule
        Send a message    & +2 XP & Per message sent \\
        Another user joins & +5 XP & Per join event received \\
        Session time       & +1 XP & Per minute of active session \\
        Streak bonus       & +10 XP & Every 5 consecutive messages \\
        \bottomrule
    \end{tabularx}
\end{table}

\subsection{Rank Ladder}

\begin{table}[H]
    \centering
    \caption{Rank Progression}
    \begin{tabularx}{\textwidth}{c l c X}
        \toprule
        \textbf{Level} & \textbf{Rank Name} & \textbf{XP Required} & \textbf{Next Rank At} \\
        \midrule
        1 & NEWBIE    & 0   & 50 XP \\
        2 & SQUIRE    & 50  & 150 XP \\
        3 & KNIGHT    & 150 & 320 XP \\
        4 & CHAMPION  & 320 & 600 XP \\
        5 & WARLORD   & 600 & 1000 XP \\
        6 & LEGEND    & 1000 & — (max rank) \\
        \bottomrule
    \end{tabularx}
\end{table}

When a user levels up, a ``LEVEL UP!'' toast notification appears with an 8-bit fanfare sound.


% ══════════════════════════════════════════════════════════════════
% 8. KEY FEATURES DEMONSTRATION
% ══════════════════════════════════════════════════════════════════
\section{Key Features — Screenshots}

This section provides visual evidence of all implemented features.

\subsection{User Join/Leave Notifications}

\begin{figure}[H]
    \centering
    \fbox{\parbox{0.85\textwidth}{\centering\vspace{3cm}
        \textbf{[PLACEHOLDER — Join/Leave Notifications Screenshot]}\\[0.5cm]
        Insert a screenshot showing join and leave system messages\\
        in the chat area (e.g., ``Alice joined the chat'', ``Bob left the chat'').
    \vspace{3cm}}}
    \caption{User Join and Leave Notifications}
    \label{fig:joinleave}
\end{figure}

\subsection{Real-Time Message Broadcasting}

\begin{figure}[H]
    \centering
    \fbox{\parbox{0.85\textwidth}{\centering\vspace{3cm}
        \textbf{[PLACEHOLDER — Multi-User Chat Screenshot]}\\[0.5cm]
        Insert a screenshot showing messages from multiple users\\
        displayed in real-time with avatars and timestamps.
    \vspace{3cm}}}
    \caption{Real-Time Message Broadcasting Across Multiple Users}
    \label{fig:broadcast}
\end{figure}

\subsection{Online Users List}

\begin{figure}[H]
    \centering
    \fbox{\parbox{0.85\textwidth}{\centering\vspace{3cm}
        \textbf{[PLACEHOLDER — Online Users Sidebar Screenshot]}\\[0.5cm]
        Insert a screenshot showing the left sidebar with\\
        multiple users listed with their avatars and the ``YOU'' tag.
    \vspace{3cm}}}
    \caption{Online Users Sidebar with Live Count}
    \label{fig:onlineusers}
\end{figure}

\subsection{Typing Indicator}

\begin{figure}[H]
    \centering
    \fbox{\parbox{0.85\textwidth}{\centering\vspace{2.5cm}
        \textbf{[PLACEHOLDER — Typing Indicator Screenshot]}\\[0.5cm]
        Insert a screenshot showing the typing indicator bar\\
        (e.g., ``ALICE TYPING...'') below the messages area.
    \vspace{2.5cm}}}
    \caption{Real-Time Typing Indicator}
    \label{fig:typing}
\end{figure}

\subsection{Emoji Picker}

\begin{figure}[H]
    \centering
    \fbox{\parbox{0.85\textwidth}{\centering\vspace{3cm}
        \textbf{[PLACEHOLDER — Emoji Picker Screenshot]}\\[0.5cm]
        Insert a screenshot showing the emoji picker panel\\
        open above the message input area.
    \vspace{3cm}}}
    \caption{Quick React Emoji Picker}
    \label{fig:emoji}
\end{figure}

\subsection{File Attachment}

\begin{figure}[H]
    \centering
    \fbox{\parbox{0.85\textwidth}{\centering\vspace{3cm}
        \textbf{[PLACEHOLDER — File Attachment Screenshot]}\\[0.5cm]
        Insert a screenshot showing a sent message with an\\
        image attachment rendered inline, or a file download card.
    \vspace{3cm}}}
    \caption{File Attachment — Image Preview and File Card}
    \label{fig:attachment}
\end{figure}

\subsection{Message Delivery Receipts}

\begin{figure}[H]
    \centering
    \fbox{\parbox{0.85\textwidth}{\centering\vspace{3cm}
        \textbf{[PLACEHOLDER — Message Receipts Screenshot]}\\[0.5cm]
        Insert a screenshot showing messages with receipt emojis\\
        (😴 for sent, 😃 for partial, 😎 for delivered to all).
    \vspace{3cm}}}
    \caption{Message Delivery Receipts}
    \label{fig:receipts}
\end{figure}

\subsection{Connection Status \& Auto-Reconnect}

\begin{figure}[H]
    \centering
    \fbox{\parbox{0.85\textwidth}{\centering\vspace{2.5cm}
        \textbf{[PLACEHOLDER — Connection Status Screenshot]}\\[0.5cm]
        Insert screenshots showing the connection status indicator\\
        in ``ONLINE'', ``OFFLINE'', and ``WAIT...'' (reconnecting) states.
    \vspace{2.5cm}}}
    \caption{Connection Status Indicator}
    \label{fig:connection}
\end{figure}

\subsection{Gamification — XP Bar \& Level Up}

\begin{figure}[H]
    \centering
    \fbox{\parbox{0.85\textwidth}{\centering\vspace{3cm}
        \textbf{[PLACEHOLDER — Gamification Screenshot]}\\[0.5cm]
        Insert a screenshot showing the XP bar, rank display,\\
        message count, streak count, and optionally the level-up toast.
    \vspace{3cm}}}
    \caption{Gamification — XP Progress Bar, Rank, and Session Stats}
    \label{fig:gamification}
\end{figure}

\subsection{Multiple Clients on Laboratory Machines}

\begin{figure}[H]
    \centering
    \fbox{\parbox{0.85\textwidth}{\centering\vspace{4cm}
        \textbf{[PLACEHOLDER — Lab Deployment Screenshot]}\\[0.5cm]
        Insert a screenshot or photo showing the application running\\
        on multiple laboratory machines simultaneously, with 4 users\\
        chatting in real-time.
    \vspace{4cm}}}
    \caption{Application Deployed on 4 Laboratory Machines}
    \label{fig:labdeploy}
\end{figure}


% ══════════════════════════════════════════════════════════════════
% 9. DEPLOYMENT & SETUP
% ══════════════════════════════════════════════════════════════════
\section{Deployment \& Setup Instructions}

\subsection{Prerequisites}
\begin{itemize}[leftmargin=*]
    \item Python 3.10 or higher installed on the server machine.
    \item \texttt{pip} package manager.
    \item All machines connected to the same local network (LAN).
    \item A modern web browser (Chrome, Firefox, Edge) on each client machine.
\end{itemize}

\subsection{Installation Steps}

\begin{lstlisting}[language=bash, caption={Installation Commands}, label={lst:install}]
# 1. Clone the repository
git clone <repository-url>
cd group-chat-app

# 2. Create the environment file
cp .env.example .env
# Edit .env to change ports if needed (default: 8000, 5000)

# 3. Install Python dependencies
cd server
pip install -r requirements.txt
cd ..
\end{lstlisting}

\subsection{Running the Application}

Two terminal windows are required on the server machine:

\begin{lstlisting}[language=bash, caption={Starting the Servers}, label={lst:start}]
# Terminal 1 — Start the Backend Server
cd server
python server.py
# Output: WebSocket: ws://0.0.0.0:8000/ws

# Terminal 2 — Start the Frontend Server
cd client
python serve.py
# Output: URL: http://localhost:5000
\end{lstlisting}

\subsection{Lab Deployment (4 Machines)}

\begin{enumerate}[leftmargin=*]
    \item \textbf{Machine 1 (Server):} Run both \texttt{server.py} and \texttt{serve.py}. Note its IP address (e.g., \texttt{192.168.1.100}).
    \item \textbf{Machines 1–4 (Clients):} Open \texttt{http://192.168.1.100:5000} in a web browser.
    \item Each user enters a unique username, selects an avatar, and clicks ``Start Quest''.
    \item The WebSocket client automatically connects to \texttt{ws://192.168.1.100:8000/ws} (uses \texttt{window.location.hostname}).
\end{enumerate}

\subsection{Environment Configuration}

\begin{lstlisting}[language=bash, caption={.env Configuration File}, label={lst:env}]
# Backend (FastAPI WebSocket server)
BACKEND_PORT=8000

# Frontend (static file server)
FRONTEND_PORT=5000
\end{lstlisting}

The \texttt{.env} file is loaded by both server scripts via \texttt{python-dotenv}. It is excluded from version control via \texttt{.gitignore}.


% ══════════════════════════════════════════════════════════════════
% 10. CORS & SECURITY
% ══════════════════════════════════════════════════════════════════
\section{CORS \& Security Considerations}

\subsection{CORS Configuration}
The backend server enables Cross-Origin Resource Sharing (CORS) to allow the frontend (running on a different port) to communicate with the backend:

\begin{lstlisting}[language=Python, caption={CORS Middleware Configuration}, label={lst:cors}]
app.add_middleware(
    CORSMiddleware,
    allow_origins=["*"],
    allow_credentials=True,
    allow_methods=["*"],
    allow_headers=["*"],
)
\end{lstlisting}

\textit{Note:} The wildcard (\texttt{*}) origin is used for simplicity in the lab environment. In a production deployment, this should be restricted to specific allowed origins.

\subsection{Security Measures}
\begin{itemize}[leftmargin=*]
    \item \textbf{XSS Prevention:} All user-generated text is sanitized using \texttt{escapeHtml()} before rendering in the DOM.
    \item \textbf{Input Validation:} The server validates usernames (non-empty, unique) and enforces the join protocol (first message must be type \texttt{join}).
    \item \textbf{File Upload Safety:} Uploaded files are renamed with UUID-generated filenames to prevent path traversal and naming conflicts.
    \item \textbf{Connection Cleanup:} Broken WebSocket connections are detected and cleaned up automatically during broadcast operations.
\end{itemize}


% ══════════════════════════════════════════════════════════════════
% 11. TESTING
% ══════════════════════════════════════════════════════════════════
\section{Testing}

The application was tested across multiple scenarios to ensure reliability and correctness.

\subsection{Test Cases}

\begin{longtable}{p{0.5cm} p{5cm} p{4cm} p{3cm}}
    \toprule
    \textbf{\#} & \textbf{Test Case} & \textbf{Expected Result} & \textbf{Status} \\
    \midrule
    \endhead
    1 & Single user joins with a valid username & Welcome message displayed, user appears in sidebar & \textbf{Pass} \\
    \midrule
    2 & Two users join with same username & Second user receives ``Username taken'' error & \textbf{Pass} \\
    \midrule
    3 & User joins with empty username & ``Enter a name first'' error displayed & \textbf{Pass} \\
    \midrule
    4 & User sends a text message & Message appears for all connected users in real-time & \textbf{Pass} \\
    \midrule
    5 & User sends a message with image attachment & Image rendered inline for all users & \textbf{Pass} \\
    \midrule
    6 & User disconnects (closes tab) & Leave notification broadcast, user removed from sidebar & \textbf{Pass} \\
    \midrule
    7 & Server goes down while client is connected & Client shows ``OFFLINE'', auto-reconnects when server restarts & \textbf{Pass} \\
    \midrule
    8 & New user joins ongoing chat & Receives last 50 messages as history & \textbf{Pass} \\
    \midrule
    9 & User types a message & ``USERNAME TYPING'' indicator shown to others & \textbf{Pass} \\
    \midrule
    10 & All users leave and one rejoins within 60s & Message history is preserved & \textbf{Pass} \\
    \midrule
    11 & All users leave, no one rejoins for 60s & Message history is cleared & \textbf{Pass} \\
    \midrule
    12 & User clicks ``Leave'' button & Graceful disconnect, returns to login screen & \textbf{Pass} \\
    \midrule
    13 & 4 users chat simultaneously on LAN & All messages delivered in real-time to all clients & \textbf{Pass} \\
    \midrule
    14 & Message delivery receipt: all online & 😎 emoji displayed on sender's message & \textbf{Pass} \\
    \midrule
    15 & Message delivery receipt: alone in room & 😴 emoji displayed on sender's message & \textbf{Pass} \\
    \bottomrule
\end{longtable}

\subsection{Testing Environment}

\begin{figure}[H]
    \centering
    \fbox{\parbox{0.85\textwidth}{\centering\vspace{3cm}
        \textbf{[PLACEHOLDER — Testing Screenshot]}\\[0.5cm]
        Insert a screenshot showing the backend server terminal\\
        logs during a test session (user joins, messages, leaves).
    \vspace{3cm}}}
    \caption{Backend Server Terminal Logs During Testing}
    \label{fig:testing}
\end{figure}


% ══════════════════════════════════════════════════════════════════
% 12. CHALLENGES & SOLUTIONS
% ══════════════════════════════════════════════════════════════════
\section{Challenges \& Solutions}

\begin{table}[H]
    \centering
    \caption{Challenges Encountered and Solutions Applied}
    \begin{tabularx}{\textwidth}{X X}
        \toprule
        \textbf{Challenge} & \textbf{Solution} \\
        \midrule
        Handling broken WebSocket connections during broadcast &
        Wrapped each \texttt{send\_json()} in a try-except block; broken connections are collected and cleaned up after the broadcast loop. \\
        \midrule
        Preventing duplicate message display (sender sees own message twice) &
        Implemented optimistic rendering on the client side; server echo for the sender's own username is suppressed. \\
        \midrule
        Cross-origin requests between frontend (port 5000) and backend (port 8000) &
        Configured CORS middleware on the backend to allow all origins. \\
        \midrule
        File attachment URLs not resolving on other clients &
        Absolutized file URLs on the client side by prepending the backend host and port to relative paths. \\
        \midrule
        Stale message history persisting across different sessions &
        Implemented a 60-second cleanup timer that clears history when the room is empty. Timer is cancelled if someone rejoins. \\
        \midrule
        Network instability causing permanent disconnections &
        Implemented auto-reconnect with exponential backoff (1s to 10s max delay). \\
        \bottomrule
    \end{tabularx}
\end{table}


% ══════════════════════════════════════════════════════════════════
% 13. WORKFLOW SUMMARY
% ══════════════════════════════════════════════════════════════════
\section{Workflow Summary}

The following is a step-by-step summary of the complete application workflow:

\begin{enumerate}[leftmargin=*]
    \item \textbf{Server Startup:} The administrator starts the backend WebSocket server (\texttt{server.py}) and the frontend static file server (\texttt{serve.py}) on the designated host machine.

    \item \textbf{Client Access:} Each user navigates to \texttt{http://<server-ip>:5000} in their browser. The browser downloads the HTML, CSS, and JavaScript files.

    \item \textbf{Login:} The user selects an avatar from the 12 available options, enters a unique username (max 20 characters), and clicks ``START QUEST''.

    \item \textbf{WebSocket Handshake:} The JavaScript client initiates a WebSocket connection to \texttt{ws://<server-ip>:8000/ws}. Upon successful TCP handshake and HTTP upgrade, the connection is established.

    \item \textbf{Join Protocol:} The client sends a \texttt{join} message. The server validates the username, registers the user, sends a welcome message, broadcasts a join notification, and syncs the user list and message history.

    \item \textbf{Real-Time Chat:} Users exchange messages in real-time. Each message is broadcast to all connected clients. The sender sees instant feedback via optimistic rendering and delivery receipts.

    \item \textbf{Interactive Features:} Users can send emojis, attach files, and see typing indicators — all communicated over the same WebSocket connection.

    \item \textbf{Gamification:} As users interact, they earn XP, build streaks, and progress through ranks — displayed on the sidebar.

    \item \textbf{Disconnect:} When a user leaves (intentionally or due to network issues), the server notifies remaining users, updates the user list, and handles cleanup. The client auto-reconnects on unintentional disconnections.

    \item \textbf{Session End:} When all users have left, the server waits 60 seconds before clearing the message history, preparing for the next session.
\end{enumerate}


% ══════════════════════════════════════════════════════════════════
% 14. FUTURE ENHANCEMENTS
% ══════════════════════════════════════════════════════════════════
\section{Future Enhancements}

While the current implementation fulfills all assignment requirements and includes several bonus features, the following enhancements could be considered for future development:

\begin{itemize}[leftmargin=*]
    \item \textbf{Multiple Chat Rooms:} Support for creating and joining different chat rooms/channels.
    \item \textbf{Private Messaging:} Direct one-to-one messaging between users.
    \item \textbf{Persistent Storage:} Database integration (e.g., SQLite, PostgreSQL) for message history and user accounts.
    \item \textbf{User Authentication:} Password-protected user accounts with login sessions.
    \item \textbf{End-to-End Encryption:} Encrypting messages so only intended recipients can read them.
    \item \textbf{Message Reactions:} Allowing users to react to specific messages with emojis.
    \item \textbf{Read Receipts:} Tracking which users have actually read a message (not just received it).
    \item \textbf{HTTPS/WSS:} TLS encryption for secure communication in production environments.
    \item \textbf{Mobile Responsive Design:} Optimizing the UI for mobile devices and tablets.
\end{itemize}


% ══════════════════════════════════════════════════════════════════
% 15. CONCLUSION
% ══════════════════════════════════════════════════════════════════
\section{Conclusion}

This project successfully demonstrates the implementation of a \textbf{real-time Group Chat Application using WebSockets}. The application meets all the minimum requirements specified in the tutorial assignment:

\begin{itemize}[leftmargin=*]
    \item[\checkmark] User join/leave notifications.
    \item[\checkmark] Real-time message broadcasting to all connected users.
    \item[\checkmark] Usernames for identifying participants.
    \item[\checkmark] Graceful handling of client disconnections.
\end{itemize}

Beyond the requirements, the application includes several advanced features — delivery receipts, file attachments, typing indicators, auto-reconnect, message history, an emoji picker, and a complete gamification system — all wrapped in a visually engaging 8-bit pixel-art themed interface.

The project demonstrates practical understanding of key computer networking concepts: persistent TCP connections (WebSocket protocol), client-server architecture, message serialization (JSON), asynchronous I/O (Python asyncio with FastAPI), and cross-origin resource sharing (CORS).

The application was successfully deployed and tested on laboratory machines with 4 simultaneous users communicating in real-time.


% ══════════════════════════════════════════════════════════════════
% 16. REFERENCES
% ══════════════════════════════════════════════════════════════════
\section{References}

\begin{enumerate}[leftmargin=*]
    \item RFC 6455 — The WebSocket Protocol. \url{https://datatracker.ietf.org/doc/html/rfc6455}
    \item FastAPI Documentation — WebSockets. \url{https://fastapi.tiangolo.com/advanced/websockets/}
    \item MDN Web Docs — WebSocket API. \url{https://developer.mozilla.org/en-US/docs/Web/API/WebSocket}
    \item Uvicorn — ASGI Server. \url{https://www.uvicorn.org/}
    \item FastAPI Official Documentation. \url{https://fastapi.tiangolo.com/}
    \item Python \texttt{asyncio} Documentation. \url{https://docs.python.org/3/library/asyncio.html}
    \item MDN Web Docs — Web Audio API. \url{https://developer.mozilla.org/en-US/docs/Web/API/Web_Audio_API}
\end{enumerate}


% ══════════════════════════════════════════════════════════════════
% APPENDIX
% ══════════════════════════════════════════════════════════════════
\newpage
\appendix
\section{Appendix: Client URL for TA Testing}

% ── UPDATE this with your actual deployment URL ──
\begin{center}
    \large
    \textbf{Client URL:} \texttt{http://<YOUR-MACHINE-IP>:5000}
\end{center}

\textit{Note:} Ensure the backend server (\texttt{server.py}) and frontend server (\texttt{serve.py}) are both running on the specified machine before accessing this URL.


\section{Appendix: Source Code Repository}

% ── UPDATE this with your actual repository URL ──
The complete source code is available at:
\begin{center}
    \url{https://github.com/<your-username>/group-chat-app}
\end{center}


\end{document}
